\chapter*{Conclusion générale}
\addcontentsline{toc}{chapter}{Conclusion générale}
L'analyse combine deux sources officielles qui remplissent des rôles distincts. L'IPC établit la chronologie des variations de prix. L'ENNVM renseigne sur la répartition des dépenses des ménages selon leur milieu et leur quintile de dépenses. La période 2022--2023 concentre les pics de variation de l'IPC général dans l'échantillon. L'alimentation contribue à caractériser cet épisode par un glissement particulièrement élevé en février 2023.

Dans le modèle à poids fixes, la variation du coût et sa contribution au budget initial dépendent du périmètre. Pour l'ensemble national, de septembre 2022 à octobre 2026, la hausse historique est plus élevée dans le sous-panier prudent, tandis que la contribution au budget initial est plus forte dans le périmètre élargi. Le choc uniforme confirme que le taux du sous-panier est identique pour tous les profils et que le supplément dépend de la couverture. Le choc alimentaire montre que la charge marginale dépend de la part initiale du poste alimentaire. Ces résultats portent sur l'exposition des budgets sous les hypothèses choisies. Des revenus comparables, des quantités achetées et des prix locaux seraient nécessaires pour prolonger l'analyse du pouvoir d'achat et de la consommation.

L'introduction de ressources nominales hypothétiques précise le seuil de compensation. Une hausse de 5\,\% suffit exactement face au choc uniforme de 5\,\% dans les deux périmètres. Après le choc combiné, une hausse de 10\,\% suffit pour 14 profils dans le périmètre élargi et pour aucun dans le prudent. Ces résultats décrivent le nombre d'unités accessibles du sous-panier et dépendent des prix, du périmètre et des ressources supposées.

L'analyse de robustesse montre que le supplément alimentaire marginal reste supérieur pour Rural C1 à celui d'Urbain C5 sous les variantes de poids testées et sous l'hypothèse de prix communs. Le classement des périmètres pour la variation du coût historique change selon la fenêtre. Le dépassement de 10\,\% pour Rural C1 après le choc combiné dans l'élargi dépend aussi de la pondération. Ces constats circonscrivent les conclusions aux périodes, profils et variantes étudiés.

Les perspectives consistent à rapprocher plus finement les nomenclatures, à recueillir des prix ruraux comparables et à disposer de mesures répétées des revenus, dépenses et quantités achetées. Ces données permettraient de confronter les taux nominaux hypothétiques aux évolutions effectivement observées, de tester la stabilité des poids et d'étudier les substitutions ainsi que l'hétérogénéité au sein de chaque profil.
